\documentclass{article}
\usepackage{graphicx}
\usepackage{amsmath}
\usepackage[T2A]{fontenc}  % For Cyrillic fonts
\usepackage[utf8]{inputenc}  % For UTF-8 encoding
\usepackage[russian]{babel}  % For Russian language support

\title{Algorithms 1 \\ Homework 1}
\author{Nikita Zagainov}
\date{September 2024}

\begin{document}

\maketitle
\section{Asymptotic complexity}
\begin{enumerate}
    \item
          \begin{enumerate}
              \item[1.]
                  (a) Да, возможно. При
                  \begin{align*}
                       & f(n) = \Theta(n^2),              \\
                       & g(n) = \Theta(\frac{n}{\log{n}})
                  \end{align*}
                  выполняется
                  \[
                      h(n) = \Theta{(n\log{n})}
                  \]
                  (b) Нет, невозможно. Для увеличения асимптотической оценки $h(n)$ нужно взять максимально быстрорастущий $f(n) = \Theta(n^2)$ и минимально быстрорастущийй $g(n) = \Theta(1)$, но это даёт только $h(n) = \Theta(n^2)$, и это максимальное значение.
              \item[2.]
                  По аналогичным рассуждениям с 1.1.b, получем следующюю верхнюю оценку:
                  \begin{align*}
                       & h(n) = \Theta(n^2), \ \text{if} \ f(n) = \Theta(n^2), \ g(n) = \Theta(1)
                  \end{align*}
                  Нижней асимптотической границы у $h(n)$, так как её нет у функции $f(n)$.
          \end{enumerate}
    \item
          \[
              \sqrt{i^3 + 2i + 5} = i^{\frac{3}{2}}\Theta(1 + \frac{1}{i})
          \]
          Следовательно,
          \[
              \sum_{i=1}^{n}{\sqrt{i^3 + 2i + 5}} = \Theta{(\int_0^n{i^{\frac{3}{2}}di})} = \Theta(\frac{2}{5}x^{\frac{5}{2}}) = \Theta(x^{\frac{5}{2}})
          \]
    \item
          Мы рассматриваем случаи, когда $\alpha > 0$, то есть, когда $\sum_{i=1}^{n}{i^{\alpha}} \to \infty, \ \text{if} \ i \to \infty$
          \begin{align*}
               & \sum_{i=1}^{n}{i^{\alpha}} = \Theta(\int_0^n{i^{\alpha} di}) = \Theta(\frac{1}{1 + \alpha}n^{\alpha + 1}) = \Theta(n^{\alpha + 1})
          \end{align*}

    \item
          \begin{align*}
               & \sum_{k=1}^{n}{\frac{1}{k}} = \Theta(\int_1^n{\frac{1}{k} dk}) = \Theta(\log{n})
          \end{align*}
\end{enumerate}

\end{document}
